% Part: lambda-calculus
% Chapter: introduction
% Section: basic-pr-lambda

\documentclass[../../../include/open-logic-section]{subfiles}

\begin{document}

\olfileid{lam}{int}{bas}
\olsection{الدوال العودية البدائية الأساسية \usetoken{S}{lambda definable}}

\begin{lem}
كل واحدة من الدوال~$\Zero$ و$\Succ$ و$\Proj{\OLId{latin-ordinary}{n}}{\OLId{latin-ordinary}{i}}$ هي !!{lambda definable}.
\end{lem}

\begin{proof}
أما~$\Zero$ فيمثلها الحد
$\lambd[\OLId{latin-ordinary}{x}][\lambd[\OLId{latin-ordinary}{y}][\OLId{latin-ordinary}{y}]]$.

وأما دالة الخلف~$\Succ$ فتعريفها
$\fn{Succ}(\OLId{latin-ordinary}{u}) = \lambd[\OLId{latin-ordinary}{x}][\lambd[\OLId{latin-ordinary}{y}][\OLId{latin-ordinary}{x}(uxy)]]$. ولتتأمل وجه هذا التعريف:
لكل عدد~$\num{\OLId{latin-ordinary}{n}}$ نعدّه مكرّرًا، ولكل دالة~$\OLId{latin-ordinary}{f}$، تكون
$\fn{Succ}(\num{\OLId{latin-ordinary}{n}},\OLId{latin-ordinary}{f})$ دالة تأخذ~$\OLId{latin-ordinary}{y}$، فتطبق~$\OLId{latin-ordinary}{f}$
$\OLId{latin-ordinary}{n}$ مرة، بادئةً من~$\OLId{latin-ordinary}{y}$، ثم تزيد تطبيقًا واحدًا.

وأما الإسقاطات فأمرها بيّن:
$\fn{Proj}^\OLId{latin-ordinary}{n}_\OLId{latin-ordinary}{i}(\OLId{latin-ordinary}{x}_\OLMathNumeral{0}, \dots, \OLId{latin-ordinary}{x}_{\OLId{latin-ordinary}{n}-\OLMathNumeral{1}}) = \OLId{latin-ordinary}{x}_\OLId{latin-ordinary}{i}$. فمقتضى اصطلاحنا أن
$\fn{Proj}^\OLId{latin-ordinary}{n}_\OLId{latin-ordinary}{i}$ هي حد لامبدا
$\lambd[\OLId{latin-ordinary}{x}_\OLMathNumeral{0}][\dots \lambd[\OLId{latin-ordinary}{x}_{\OLId{latin-ordinary}{n}-\OLMathNumeral{1}}][\OLId{latin-ordinary}{x}_\OLId{latin-ordinary}{i}]]$.
\end{proof}

\end{document}

